% Introducción — síntesis del Cap. I + estructura del documento (Norma §6.1). Actualizar al cerrar el Cap. IV.
\section*{INTRODUCCIÓN}
\addcontentsline{toc}{section}{INTRODUCCIÓN}

ISI Mustang Bolivia ejecuta proyectos de ingeniería, procura y construcción (EPC) para los sectores de gas y petróleo, energía y minería. Durante más de una década planificó y controló esos proyectos con un sistema propio, el Portal ISIMustang Bolivia, que acumuló el registro de sus presupuestos, costos, horas y certificaciones, y hoy lo reemplaza por un ERP nuevo de desarrollo propio. Ese histórico se usa solo para saber qué ocurrió: las desviaciones de costo y de plazo se detectan cuando ya se materializaron y su corrección es más costosa.

El presente proyecto de grado desarrolla un módulo de inteligencia predictiva integrado al ERP de la empresa, que aplica técnicas de aprendizaje automático sobre ese histórico para anticipar desviaciones en tres casos de uso: la desviación en certificaciones, el sobrecosto de proyecto y el retraso en la fecha de cierre. El módulo predictivo es el eje del trabajo. La modernización del ERP que el trabajo también aborda es la condición para integrarlo y usarlo: la reimplementación de los procesos del portal anterior que producen los datos que el módulo necesita, la autenticación centralizada con Keycloak, que da una identidad común a todos los servicios, los reportes descriptivos en Power BI y un servidor MCP que expone los datos del ERP y las predicciones a clientes de inteligencia artificial.

El trabajo corresponde a una investigación aplicada y se desarrolla con una metodología iterativa e incremental en cinco incrementos. Parte del análisis de los datos del portal anterior, que determina qué puede predecirse, con qué variables y sobre cuántos casos; esa evidencia orienta el diseño de los modelos y de la arquitectura, y cada incremento se construye y valida sobre el ERP de la empresa.

El documento se organiza en cuatro capítulos. El Capítulo I, Marco General, presenta los antecedentes, el problema, los objetivos, la justificación, la metodología y la delimitación del trabajo. El Capítulo II, Marco Teórico y Tecnológico, desarrolla los fundamentos de los sistemas ERP, el aprendizaje automático con muestras pequeñas, la ingeniería de datos, la analítica predictiva en la gestión de proyectos EPC, la ingeniería de software, las tecnologías del sistema, la gestión de identidad, el protocolo MCP y la inteligencia de negocios. El Capítulo III, Análisis y Diseño del Sistema Propuesto, analiza los datos históricos y los procesos del portal anterior, define los requerimientos del sistema y diseña su arquitectura. El Capítulo IV, Construcción e Implementación del Sistema, documenta los cinco incrementos y la evaluación de los modelos y del sistema con usuarios clave. Cierran el documento las conclusiones y recomendaciones, las referencias y los apéndices.
